\documentclass[11pt,a4paper]{report}

\usepackage[utf8]{inputenc}
\usepackage[T1]{fontenc}
\usepackage[a4paper,margin=2.5cm]{geometry}
\usepackage{longtable}
\usepackage{booktabs}
\usepackage{array}
\usepackage{tabularx}
\usepackage{ltablex}
\usepackage{enumitem}
\usepackage{xcolor}
\usepackage{hyperref}
\usepackage{listings}
\usepackage[most]{tcolorbox}
\usepackage{titlesec}
\usepackage{setspace}
\usepackage{parskip}

\hypersetup{
  colorlinks=true,
  linkcolor=blue!50!black,
  urlcolor=blue!50!black
}

\definecolor{codebg}{RGB}{247,247,247}
\definecolor{notebg}{RGB}{245,249,255}
\definecolor{noteedge}{RGB}{72,110,170}

\lstdefinestyle{mgpy}{
  basicstyle=\ttfamily\small,
  backgroundcolor=\color{codebg},
  frame=single,
  framerule=0pt,
  breaklines=true,
  columns=fullflexible,
  keepspaces=true
}

\newtcolorbox{guidenote}[1][]{
  colback=notebg,
  colframe=noteedge,
  boxrule=0.6pt,
  arc=2pt,
  left=8pt,
  right=8pt,
  top=6pt,
  bottom=6pt,
  title=#1
}

\newcommand{\code}[1]{\texttt{#1}}
\newcolumntype{Y}{>{\raggedright\arraybackslash}X}

\titleformat{\chapter}[display]
  {\normalfont\bfseries}
  {\chaptertitlename\ \thechapter}{10pt}{\Large}

\title{MicroGridsPy Planning User Guide}
\author{MicroGridsPy Planning}
\date{\today}

\begin{document}

\maketitle
\tableofcontents

\part{Shared Platform Workflow}

\chapter{Introduction}\label{sec:introduction}

MicroGridsPy Planning is the planning module of the MicroGridsPy ecosystem. It is implemented as a Streamlit application backed by Linopy optimization and a project-folder workflow based on JSON, YAML, and CSV files.

The application supports two complementary planning modes:

\begin{itemize}[leftmargin=1.5em]
  \item \textbf{Multi-Year}: an explicit dynamic formulation with year-by-year evolution, investment-step logic, and cohort-style capacity tracking over the planning horizon.
  \item \textbf{Typical-Year}: a reduced steady-state formulation in which one representative operating year is repeated and investment decisions are evaluated in annuity form.
\end{itemize}

This guide adopts the Multi-Year workflow as the \emph{reference} implementation. Shared workflow, interface, directory, input, and conditional-feature logic are therefore defined once in the Multi-Year chapters. The Typical-Year chapters later describe only the simplifications and differences required when the reduced formulation is selected.

\begin{guidenote}[Reference framing]
In the current implementation, Multi-Year is the general planning mode and Typical-Year is its reduced special case. When a concept exists in both modes, this guide defines it in the Multi-Year chapters first and then refers back to it from the Typical-Year chapters.
\end{guidenote}

\chapter{General Workflow}\label{sec:general-workflow}

The minimum workflow for any project is:

\begin{enumerate}[leftmargin=1.5em]
  \item Launch the Streamlit application and open \textbf{Project Setup}.
  \item Create a project folder under \path{projects/<project\_name>/}.
  \item Select the planning formulation.
  \item Configure the high-level project settings that determine which templates are generated.
  \item Click \textbf{Initialize project and generate templates}.
  \item Edit the generated input files under \path{projects/<project\_name>/inputs/}.
  \item Open \textbf{Data Audit and Visualization} to confirm that the dataset loads and the dimensions match the intended study.
  \item Open \textbf{Model Optimization} and solve the model.
  \item Open \textbf{Results} to inspect sizing, operation, costs, emissions, and exported files.
\end{enumerate}

This workflow is shared by both planning modes. The main differences lie in how the formulation is initialized and in how years, investment steps, and time-series inputs are organized.

\chapter{Streamlit Interface}\label{sec:streamlit-interface}

\section{Home}

\path{Home.py} is the landing page. It links to:

\begin{itemize}[leftmargin=1.5em]
  \item \path{pages/0\_Project\_Setup.py}
  \item \path{pages/1\_Data\_Audit\_and\_Visualization.py}
  \item \path{pages/3\_Optimization.py}
  \item \path{pages/4\_Results.py}
\end{itemize}

The Home page introduces the two formulations at a high level but does not configure the model directly.

\section{Project Setup}

This page creates or loads projects.

\subsection{Create tab}

The Create tab controls:

\begin{itemize}[leftmargin=1.5em]
  \item project name and description,
  \item formulation selection,
  \item grid and export settings,
  \item renewable-resource count and labels,
  \item uncertainty settings,
  \item battery and generator advanced-model options,
  \item optimization constraints,
  \item template generation.
\end{itemize}

\subsection{Load tab}

The Load tab marks an existing project as active without regenerating inputs.

\begin{guidenote}[Important behavior]
If a project is created with a name that already exists, the generated files in \path{inputs/} are overwritten using the current Project Setup selections. Existing user-edited inputs should therefore be backed up or reviewed before regeneration.
\end{guidenote}

\section{Data Audit and Visualization}

This page reads \path{inputs/formulation.json}, initializes the formulation-specific sets, loads the canonical project dataset, and displays:

\begin{itemize}[leftmargin=1.5em]
  \item required and optional input-file status,
  \item dataset-loading status,
  \item coordinate summaries,
  \item parameter summaries,
  \item optimization-constraint summaries,
  \item advanced-curve diagnostics,
  \item grid-availability controls for on-grid projects,
  \item time-series plots.
\end{itemize}

For on-grid projects, \path{grid\_availability.csv} is treated as a derived artifact and can be regenerated from \path{grid.yaml}.

\section{Model Optimization}

This page:

\begin{itemize}[leftmargin=1.5em]
  \item requires an active project,
  \item reads \path{inputs/formulation.json},
  \item builds the formulation-specific optimization model,
  \item lets the user select \code{highs} or \code{gurobi},
  \item optionally writes a problem file,
  \item stores a solver log under \path{projects/<project\_name>/logs/},
  \item solves the single-objective planning problem.
\end{itemize}

\section{Results}

This page dispatches to the formulation-specific results renderer. Depending on the active mode, it displays sizing, key performance indicators, operational summaries, economic breakdowns, and exported CSV or Excel artifacts.

\part{Multi-Year Planning Mode}

\chapter{Multi-Year Overview}\label{sec:multi-year-overview}

Multi-Year is implemented as \code{core\_formulation = "dynamic"}. It is the general planning workflow in the current codebase.

Compared with Typical-Year, Multi-Year is structurally richer because it makes the following dimensions explicit:

\begin{itemize}[leftmargin=1.5em]
  \item \code{year}: modeled years over the project horizon,
  \item \code{inv\_step}: investment steps and commissioning cohorts,
  \item \code{scenario}: uncertainty branches,
  \item \code{resource}: renewable-resource labels,
  \item \code{period}: hourly periods within the representative operating year.
\end{itemize}

The Multi-Year formulation should be used when the study requires explicit year evolution, time-dependent fuel prices or grid conditions, investment sequencing, battery degradation tracking, or cohort-aware capacity accounting.

\chapter{Multi-Year Initialization Logic}\label{sec:multi-year-init}

\section{Formulation switch}

Project Setup activates Multi-Year by writing
\code{core\_formulation = "dynamic"} to \path{inputs/formulation.json}.

The Multi-Year set initializer then derives:

\begin{itemize}[leftmargin=1.5em]
  \item \code{period = 0,\dots,8759},
  \item \code{year = [start\_year\_label, start\_year\_label + 1, \dots]},
  \item \code{scenario = scenario labels},
  \item \code{resource = renewable-resource labels},
  \item \code{inv\_step = [1]} if capacity expansion is disabled,
  \item \code{inv\_step = [1,2,\dots,N]} if capacity expansion is enabled.
\end{itemize}

\section{Year labels}

The current implementation requires:

\begin{itemize}[leftmargin=1.5em]
  \item \code{start\_year\_label} to be integer-like,
  \item \code{time\_horizon\_years} to be a positive integer.
\end{itemize}

If \code{start\_year\_label = 2026} and \code{time\_horizon\_years = 10}, the modeled years are:
\code{2026, 2027, \dots, 2035}.

Those labels are then reused consistently in CSV headers, yearly fuel-price lists, grid-connection timing, and yearly result exports.

\section{Capacity expansion and investment steps}

If \code{capacity\_expansion = false}, there is one investment step only and generated YAML templates use a single \code{"1"} block under \code{investment.by\_step}.

If \code{capacity\_expansion = true}, the UI requires \code{investment\_steps\_years}, for example \code{[5,5]}. The initializer then constructs:

\begin{itemize}[leftmargin=1.5em]
  \item \code{inv\_step = [1,2]},
  \item \code{inv\_step\_start\_year = [2026,2031]},
  \item \code{inv\_step\_end\_year = [2030,2035]},
  \item \code{inv\_step\_len\_years = [5,5]},
  \item \code{year\_inv\_step[year]},
  \item \code{inv\_active\_in\_year[inv\_step, year]}.
\end{itemize}

\begin{guidenote}[Capacity expansion and cohorts]
In the current backend, investment-side parameters may vary by \code{inv\_step}, but technical parameters remain shared across steps. Capacity commissioned in a step becomes a cohort or vintage tracked through later years, while the underlying technology definition stays shared.
\end{guidenote}

\section{Canonical step-key convention}

The current user-facing Multi-Year templates use canonical numeric step labels:

\begin{itemize}[leftmargin=1.5em]
  \item \code{"1"} for the single-step case,
  \item \code{"1"}, \code{"2"}, \dots for multi-step cases.
\end{itemize}

The loader still accepts lightweight legacy aliases such as \code{step\_1} and, for single-step projects only, \code{base}. New inputs should follow the numeric convention directly.

\section{Scenarios and scenario weights}

If multi-scenario mode is disabled, the application uses a single scenario, typically \code{scenario\_1}. If it is enabled, \code{multi\_scenario.scenario\_labels} defines the ordered labels and \code{multi\_scenario.scenario\_weights} defines the scenario weights. All scenario-dependent CSV headers and YAML mappings must use the same labels.

\chapter{Multi-Year Project Directory}\label{sec:multi-year-directory}

The representative Multi-Year project structure is:

\begin{lstlisting}[style=mgpy]
projects/
  <project_name>/
    inputs/
      formulation.json
      README_inputs.md
      load_demand.csv
      resource_availability.csv
      renewables.yaml
      battery.yaml
      generator.yaml
      grid.yaml                       # only if on-grid
      grid_import_price.csv           # only if on-grid
      grid_export_price.csv           # only if on-grid and export enabled
      grid_availability.csv           # derived, backend-generated
      battery_efficiency_curve.csv    # only if convex battery loss is enabled
      battery_calendar_fade_curve.csv # only if calendar fade is enabled
      generator_efficiency_curve.csv  # only if generator curve mode is enabled
    logs/
    results/
      multi_year/
        dispatch_timeseries.csv
        energy_balance.csv
        design_by_step.csv
        kpis_yearly.csv
        cashflows_discounted.csv
        scenario_costs_yearly.csv
        results_multi_year.xlsx
\end{lstlisting}

\section{Auto-generated and user-maintained files}

\begin{itemize}[leftmargin=1.5em]
  \item Auto-generated: \path{formulation.json}, \path{README\_inputs.md}, the primary input templates, solver logs, and result exports.
  \item Derived: \path{grid\_availability.csv} for on-grid projects.
  \item User-edited: the primary input templates under \path{inputs/}, except \path{grid\_availability.csv}.
\end{itemize}

\chapter{Multi-Year Input Files}\label{sec:multi-year-inputs}

\section{Master file table}

\begin{longtable}{p{0.18\textwidth}p{0.11\textwidth}p{0.12\textwidth}p{0.18\textwidth}p{0.10\textwidth}p{0.22\textwidth}}
\toprule
\textbf{File} & \textbf{Format} & \textbf{Mandatory} & \textbf{When used} & \textbf{Edited by user} & \textbf{Dimensional meaning} \\
\midrule
\endfirsthead
\toprule
\textbf{File} & \textbf{Format} & \textbf{Mandatory} & \textbf{When used} & \textbf{Edited by user} & \textbf{Dimensional meaning} \\
\midrule
\endhead
\path{formulation.json} & JSON & Yes & always & normally no & project-level flags, horizon, scenarios, constraints \\
\path{load_demand.csv} & CSV & Yes & always & Yes & hourly demand by scenario and year \\
\path{resource_availability.csv} & CSV & Yes & always & Yes & hourly renewable availability by scenario, year, and resource \\
\path{renewables.yaml} & YAML & Yes & always & Yes & renewable investment by step and shared technical data \\
\path{battery.yaml} & YAML & Yes & always & Yes & battery investment by step and shared technical and degradation data \\
\path{generator.yaml} & YAML & Yes & always & Yes & generator investment by step, shared technical data, yearly fuel prices \\
\path{grid.yaml} & YAML & Conditional & on-grid only & Yes & grid line, outages, renewable share, emissions, connection timing \\
\path{grid_import_price.csv} & CSV & Conditional & on-grid only & Yes & hourly import tariff by scenario and year \\
\path{grid_export_price.csv} & CSV & Conditional & export enabled & Yes & hourly export tariff by scenario and year \\
\path{grid_availability.csv} & CSV & Derived & on-grid only & No & backend-generated hourly availability by scenario and year \\
\path{battery_efficiency_curve.csv} & CSV & Conditional & convex battery loss & Yes & battery power-efficiency curve \\
\path{battery_calendar_fade_curve.csv} & CSV & Conditional & calendar fade enabled & Yes & calendar-ageing curve by average yearly state of charge \\
\path{generator_efficiency_curve.csv} & CSV & Conditional & generator curve mode & Yes & generator partial-load efficiency curve \\
\path{README_inputs.md} & Markdown & No & always & No & human-readable summary only \\
\bottomrule
\end{longtable}

\section{\path{inputs/formulation.json}}

This file is written by Project Setup and read by the pages, set initializer, and Multi-Year data loader.

\begin{longtable}{p{0.34\textwidth}p{0.20\textwidth}p{0.10\textwidth}p{0.10\textwidth}p{0.18\textwidth}}
\toprule
\textbf{Path} & \textbf{Meaning} & \textbf{Type} & \textbf{Unit} & \textbf{Notes} \\
\midrule
\endfirsthead
\toprule
\textbf{Path} & \textbf{Meaning} & \textbf{Type} & \textbf{Unit} & \textbf{Notes} \\
\midrule
\endhead
\code{core\_formulation} & formulation switch & string & -- & must be \code{"dynamic"} \\
\code{system\_type} & UI-level system type & string & -- & \code{"off\_grid"} or \code{"on\_grid"} \\
\code{on\_grid} & grid activation flag & bool & -- & used directly by loader logic \\
\code{grid\_allow\_export} & export activation flag & bool & -- & controls export variables and file generation \\
\code{unit\_commitment} & integer sizing toggle & bool & -- & in current code this means integer sizing, not chronological unit commitment \\
\code{start\_year\_label} & first modeled year label & int-like string or int & year & must be parseable as a positive integer \\
\code{time\_horizon\_years} & modeled horizon length & int & years & positive integer \\
\code{social\_discount\_rate} & social discount rate & number & share & used in objective accounting and reporting \\
\code{capacity\_expansion} & capacity-expansion toggle & bool & -- & activates multiple investment steps \\
\code{investment\_steps\_years} & ordered list of step durations & list[int] or null & years & required when capacity expansion is enabled; sum must equal horizon length \\
\code{multi\_scenario.enabled} & scenario replication flag & bool & -- & false means one scenario \\
\code{multi\_scenario.n\_scenarios} & number of scenarios & int & -- & should match labels and weights \\
\code{multi\_scenario.scenario\_labels} & scenario labels & list[str] & -- & must match CSV headers and YAML mappings \\
\code{multi\_scenario.scenario\_weights} & scenario weights & list[number] & share & normalized by the loader \\
\code{optimization\_constraints.enforcement} & scenario-enforcement mode & string & -- & \code{"scenario\_wise"} or \code{"expected"} \\
\code{optimization\_constraints.min\_renewable\_penetration} & renewable-share floor & number & share & typically between 0 and 1 \\
\code{optimization\_constraints.max\_lost\_load\_fraction} & unmet-load cap & number & share & typically between 0 and 1 \\
\code{optimization\_constraints.lost\_load\_cost\_per\_kwh} & unmet-load penalty & number & currency/kWh & zero disables the economic penalty \\
\code{optimization\_constraints.land\_availability\_m2} & land cap & number or null & m\textsuperscript{2} & null disables the land constraint \\
\code{optimization\_constraints.emission\_cost\_per\_kgco2e} & emissions cost & number & currency/kgCO\textsubscript{2}e & applied in objective and reporting \\
\code{system\_configuration.n\_sources} & renewable-resource count & int & -- & should match generated renewable entries \\
\code{battery\_model.loss\_model} & battery loss formulation & string & -- & \code{"constant\_efficiency"} or \code{"convex\_loss\_epigraph"} \\
\code{battery\_model.degradation\_model.cycle\_fade\_enabled} & cycle-fade toggle & bool & -- & valid in Multi-Year only with convex battery loss \\
\code{battery\_model.degradation\_model.calendar\_fade\_enabled} & calendar-fade toggle & bool & -- & valid in Multi-Year only with convex battery loss \\
\code{battery\_model.degradation\_model.*} other detailed fields & legacy detailed degradation inputs & mixed & mixed & current UI no longer writes battery-specific degradation data here; current workflow stores them in \path{battery.yaml} \\
\code{generator\_model.efficiency\_model} & generator part-load formulation & string & -- & \code{"constant\_efficiency"} or \code{"efficiency\_curve"} \\
\bottomrule
\end{longtable}

\section{\path{inputs/renewables.yaml}}

Generated schema:

\begin{lstlisting}[style=mgpy]
renewables:
  - id: res_1
    conversion_technology: Solar_PV
    resource: Solar
    investment:
      by_step:
        "1":
          nominal_capacity_kw: 1.0
          specific_investment_cost_per_kw: 0.0
          wacc: 0.0
          grant_share_of_capex: 0.0
          lifetime_years: 25
          embedded_emissions_kgco2e_per_kw: 0.0
          fixed_om_share_per_year: 0.0
          production_subsidy_per_kwh: 0.0
    technical:
      inverter_efficiency: 1.0
      specific_area_m2_per_kw: null
      max_installable_capacity_kw: null
      capacity_degradation_rate_per_year: 0.0
\end{lstlisting}

One entry is required for each renewable resource. The field \code{resource} must match the third header level of \path{resource\_availability.csv}.

\begin{longtable}{p{0.34\textwidth}p{0.19\textwidth}p{0.09\textwidth}p{0.09\textwidth}p{0.23\textwidth}}
\toprule
\textbf{Path} & \textbf{Meaning} & \textbf{Type} & \textbf{Unit} & \textbf{Behavior and notes} \\
\midrule
\endfirsthead
\toprule
\textbf{Path} & \textbf{Meaning} & \textbf{Type} & \textbf{Unit} & \textbf{Behavior and notes} \\
\midrule
\endhead
\code{renewables[].id} & internal identifier & string & -- & informational only \\
\code{renewables[].conversion\_technology} & technology label & string & -- & shared across steps and scenarios \\
\code{renewables[].resource} & resource label & string & -- & must match \path{resource\_availability.csv} and the \code{resource} coordinate \\
\code{renewables[].investment.by\_step.<step>.nominal\_capacity\_kw} & capacity per installed unit & number & kW & step-dependent \\
\code{renewables[].investment.by\_step.<step>.specific\_investment\_cost\_per\_kw} & CAPEX intensity & number & currency/kW & step-dependent; annualized with WACC and lifetime \\
\code{renewables[].investment.by\_step.<step>.wacc} & technology WACC & number & share & step-dependent \\
\code{renewables[].investment.by\_step.<step>.grant\_share\_of\_capex} & CAPEX grant share & number & share & step-dependent \\
\code{renewables[].investment.by\_step.<step>.lifetime\_years} & lifetime & number & years & step-dependent; used by replacement logic \\
\code{renewables[].investment.by\_step.<step>.embedded\_emissions\_kgco2e\_per\_kw} & embodied emissions & number & kgCO\textsubscript{2}e/kW & step-dependent \\
\code{renewables[].investment.by\_step.<step>.fixed\_om\_share\_per\_year} & fixed O\&M share & number & share/year & optional; defaults to 0.0 if omitted \\
\code{renewables[].investment.by\_step.<step>.production\_subsidy\_per\_kwh} & operating subsidy & number & currency/kWh & optional; defaults to 0.0 if omitted \\
\code{renewables[].technical.inverter\_efficiency} & conversion efficiency & number & -- & shared; usually in $(0,1]$ \\
\code{renewables[].technical.specific\_area\_m2\_per\_kw} & land coefficient & number or null & m\textsuperscript{2}/kW & shared; needed in practice when land constraint is active \\
\code{renewables[].technical.max\_installable\_capacity\_kw} & planning cap & number or null & kW & shared; null means no explicit cap \\
\code{renewables[].technical.capacity\_degradation\_rate\_per\_year} & exogenous capacity fade & number & 1/year & optional; defaults to 0.0 \\
\bottomrule
\end{longtable}

The current Multi-Year parser accepts the shared-technology schema only. Legacy branches such as top-level \code{by\_step}, \code{technical.by\_step}, and \code{operation} are no longer part of the current user workflow.

\section{\path{inputs/battery.yaml}}

Generated schema:

\begin{lstlisting}[style=mgpy]
battery:
  label: Battery
  investment:
    by_step:
      "1":
        nominal_capacity_kwh: 1.0
        specific_investment_cost_per_kwh: 0.0
        wacc: 0.0
        calendar_lifetime_years: 10
        embedded_emissions_kgco2e_per_kwh: 0.0
        fixed_om_share_per_year: 0.0
  technical:
    charge_efficiency: 0.95
    discharge_efficiency: 0.96
    initial_soc: 0.5
    depth_of_discharge: 0.8
    max_discharge_time_hours: 5.0
    max_charge_time_hours: 5.0
    max_installable_capacity_kwh: null
    efficiency_curve_csv: battery_efficiency_curve.csv
    capacity_degradation_rate_per_year: 0.0
    initial_soh: 1.0
    end_of_life_soh: 0.8
    cycle_lifetime_to_eol_cycles: 6000.0
    calendar_fade_curve_csv: battery_calendar_fade_curve.csv
    calendar_time_increment_per_year: 1.0
\end{lstlisting}

\begin{longtable}{p{0.34\textwidth}p{0.19\textwidth}p{0.09\textwidth}p{0.09\textwidth}p{0.23\textwidth}}
\toprule
\textbf{Path} & \textbf{Meaning} & \textbf{Type} & \textbf{Unit} & \textbf{Behavior and notes} \\
\midrule
\endfirsthead
\toprule
\textbf{Path} & \textbf{Meaning} & \textbf{Type} & \textbf{Unit} & \textbf{Behavior and notes} \\
\midrule
\endhead
\code{battery.label} & battery label & string & -- & reporting only \\
\code{battery.investment.by\_step.<step>.nominal\_capacity\_kwh} & energy capacity per installed unit & number & kWh & step-dependent \\
\code{battery.investment.by\_step.<step>.specific\_investment\_cost\_per\_kwh} & CAPEX intensity & number & currency/kWh & step-dependent; annualized with WACC and lifetime \\
\code{battery.investment.by\_step.<step>.wacc} & battery WACC & number & share & step-dependent \\
\code{battery.investment.by\_step.<step>.calendar\_lifetime\_years} & replacement lifetime & number & years & step-dependent \\
\code{battery.investment.by\_step.<step>.embedded\_emissions\_kgco2e\_per\_kwh} & embodied emissions & number & kgCO\textsubscript{2}e/kWh & step-dependent \\
\code{battery.investment.by\_step.<step>.fixed\_om\_share\_per\_year} & fixed O\&M share & number & share/year & optional; defaults to 0.0 if omitted \\
\code{battery.technical.charge\_efficiency} & one-way charge efficiency & number & -- & shared across steps \\
\code{battery.technical.discharge\_efficiency} & one-way discharge efficiency & number & -- & shared across steps \\
\code{battery.technical.initial\_soc} & initial state of charge & number & share & shared \\
\code{battery.technical.initial\_soh} & initial state of health & number & share & required when endogenous degradation is enabled in current generated workflow \\
\code{battery.technical.depth\_of\_discharge} & usable fraction of nominal capacity & number & share & shared \\
\code{battery.technical.max\_discharge\_time\_hours} & discharge-time proxy & number & h & shared \\
\code{battery.technical.max\_charge\_time\_hours} & charge-time proxy & number & h & shared \\
\code{battery.technical.max\_installable\_capacity\_kwh} & planning cap & number or null & kWh & shared; null means no explicit cap \\
\code{battery.technical.efficiency\_curve\_csv} & battery loss-curve filename & string or null & path & required when convex battery loss is selected \\
\code{battery.technical.capacity\_degradation\_rate\_per\_year} & exogenous capacity fade & number & 1/year & used unless calendar fade is enabled; then internally suppressed \\
\code{battery.technical.end\_of\_life\_soh} & end-of-life state of health & number & share & conditional; used for endogenous degradation \\
\code{battery.technical.cycle\_lifetime\_to\_eol\_cycles} & cycle life to end of life & number & cycles & conditional; used when cycle fade is enabled \\
\code{battery.technical.calendar\_fade\_curve\_csv} & calendar-fade filename & string & path & conditional; used when calendar fade is enabled \\
\code{battery.technical.calendar\_time\_increment\_per\_year} & calendar-ageing increment & number & increment/year & conditional; current code supports constant-per-year mode only \\
\bottomrule
\end{longtable}

\begin{guidenote}[Battery degradation ownership]
The current user-facing workflow separates high-level activation flags from battery-specific degradation data. \path{formulation.json} activates cycle fade and calendar fade, while \path{battery.yaml} stores the battery's actual degradation parameters such as \code{initial\_soh}, end-of-life state of health, cycle lifetime, and calendar-fade file paths.
\end{guidenote}

\section{\path{inputs/generator.yaml}}

Generated schema:

\begin{lstlisting}[style=mgpy]
generator:
  label: Generator
  investment:
    by_step:
      "1":
        nominal_capacity_kw: 1.0
        lifetime_years: 10
        specific_investment_cost_per_kw: 0.0
        wacc: 0.0
        embedded_emissions_kgco2e_per_kw: 0.0
        fixed_om_share_per_year: 0.0
  technical:
    nominal_efficiency_full_load: 0.30
    efficiency_curve_csv: generator_efficiency_curve.csv
    max_installable_capacity_kw: null
    capacity_degradation_rate_per_year: 0.0
fuel:
  label: Fuel
  technical:
    lhv_kwh_per_unit_fuel: 0.0
    direct_emissions_kgco2e_per_unit_fuel: 0.0
  cost:
    by_scenario:
      scenario_1:
        by_year_cost_per_unit_fuel: [0.0, 0.0]
\end{lstlisting}

\begin{longtable}{p{0.34\textwidth}p{0.19\textwidth}p{0.09\textwidth}p{0.09\textwidth}p{0.23\textwidth}}
\toprule
\textbf{Path} & \textbf{Meaning} & \textbf{Type} & \textbf{Unit} & \textbf{Behavior and notes} \\
\midrule
\endfirsthead
\toprule
\textbf{Path} & \textbf{Meaning} & \textbf{Type} & \textbf{Unit} & \textbf{Behavior and notes} \\
\midrule
\endhead
\code{generator.label} & generator label & string & -- & reporting only \\
\code{generator.investment.by\_step.<step>.nominal\_capacity\_kw} & capacity per installed unit & number & kW & step-dependent \\
\code{generator.investment.by\_step.<step>.lifetime\_years} & lifetime & number & years & step-dependent \\
\code{generator.investment.by\_step.<step>.specific\_investment\_cost\_per\_kw} & CAPEX intensity & number & currency/kW & step-dependent \\
\code{generator.investment.by\_step.<step>.wacc} & generator WACC & number & share & step-dependent \\
\code{generator.investment.by\_step.<step>.embedded\_emissions\_kgco2e\_per\_kw} & embodied emissions & number & kgCO\textsubscript{2}e/kW & step-dependent \\
\code{generator.investment.by\_step.<step>.fixed\_om\_share\_per\_year} & fixed O\&M share & number & share/year & optional; defaults to 0.0 if omitted \\
\code{generator.technical.nominal\_efficiency\_full\_load} & full-load efficiency & number & -- & shared across steps \\
\code{generator.technical.efficiency\_curve\_csv} & curve filename & string or null & path & conditional; used when generator curve mode is active \\
\code{generator.technical.max\_installable\_capacity\_kw} & planning cap & number or null & kW & shared; null means no explicit cap \\
\code{generator.technical.capacity\_degradation\_rate\_per\_year} & exogenous capacity fade & number & 1/year & optional; defaults to 0.0 \\
\code{fuel.label} & fuel label & string & -- & reporting only \\
\code{fuel.technical.lhv\_kwh\_per\_unit\_fuel} & lower heating value & number & kWh/unit fuel & shared across steps and scenarios \\
\code{fuel.technical.direct\_emissions\_kgco2e\_per\_unit\_fuel} & direct emissions factor & number & kgCO\textsubscript{2}e/unit fuel & shared across steps and scenarios \\
\code{fuel.cost.by\_scenario.<scenario>.by\_year\_cost\_per\_unit\_fuel} & yearly fuel-price trajectory & list[number] & currency/unit fuel & scenario- and year-dependent; list length must equal horizon length \\
\bottomrule
\end{longtable}

\section{\path{inputs/grid.yaml}}

This file exists only if \code{on\_grid = true}.

\begin{lstlisting}[style=mgpy]
grid:
  by_scenario:
    scenario_1:
      line:
        capacity_kw: 0.0
        transmission_efficiency: 1.0
        renewable_share: 0.0
        emissions_factor_kgco2e_per_kwh: 0.0
      outages:
        average_outages_per_year: 0.0
        average_outage_duration_minutes: 0.0
        outage_scale_od_hours: 0.6
        outage_shape_od: 0.56
        outage_seed: 0
      first_year_connection: 2026
\end{lstlisting}

\begin{longtable}{p{0.34\textwidth}p{0.19\textwidth}p{0.09\textwidth}p{0.09\textwidth}p{0.23\textwidth}}
\toprule
\textbf{Path} & \textbf{Meaning} & \textbf{Type} & \textbf{Unit} & \textbf{Behavior and notes} \\
\midrule
\endfirsthead
\toprule
\textbf{Path} & \textbf{Meaning} & \textbf{Type} & \textbf{Unit} & \textbf{Behavior and notes} \\
\midrule
\endhead
\code{grid.by\_scenario.<scenario>.line.capacity\_kw} & interconnection capacity & number & kW & scenario-dependent; same limit used for import and export \\
\code{grid.by\_scenario.<scenario>.line.transmission\_efficiency} & delivery efficiency & number & -- & scenario-dependent; must stay in $[0,1]$ \\
\code{grid.by\_scenario.<scenario>.line.renewable\_share} & renewable share of imported power & number & share & scenario-dependent; must stay in $[0,1]$ \\
\code{grid.by\_scenario.<scenario>.line.emissions\_factor\_kgco2e\_per\_kwh} & grid import emissions factor & number & kgCO\textsubscript{2}e/kWh & scenario-dependent; non-negative \\
\code{grid.by\_scenario.<scenario>.outages.average\_outages\_per\_year} & outage frequency & number & events/year & used by outage simulator \\
\code{grid.by\_scenario.<scenario>.outages.average\_outage\_duration\_minutes} & outage duration & number & minutes & used by outage simulator \\
\code{grid.by\_scenario.<scenario>.outages.outage\_scale\_od\_hours} & Weibull scale parameter & number & h & must be strictly positive \\
\code{grid.by\_scenario.<scenario>.outages.outage\_shape\_od} & Weibull shape parameter & number & -- & must be strictly positive \\
\code{grid.by\_scenario.<scenario>.outages.outage\_seed} & outage random seed & int-like number & seed & used for deterministic regeneration \\
\code{grid.by\_scenario.<scenario>.first\_year\_connection} & first connected modeled year & year label, int, null, or blank & year & null or blank means connected from horizon start \\
\bottomrule
\end{longtable}

\begin{guidenote}[Connection timing]
The preferred input for \code{first\_year\_connection} is an exact label from the modeled \code{year} coordinate, for example \code{2031}. If it is after the last modeled year, grid availability remains zero for the whole horizon.
\end{guidenote}

\section{Informational files}

\path{README\_inputs.md} is generated for human guidance only and is not parsed as a primary model input. YAML \code{meta} blocks, where present, are descriptive rather than structural.

\chapter{Multi-Year Time-Series Structure}\label{sec:multi-year-ts}

\section{\path{load_demand.csv}}

This file uses a two-row header:

\begin{itemize}[leftmargin=1.5em]
  \item row 1: scenario,
  \item row 2: year.
\end{itemize}

The first column must be \code{meta/hour}. All scenario-year combinations must exist. The file must contain exactly 8760 rows and \code{meta/hour} must run from 0 to 8759. Values are hourly demand in kWh.

\begin{lstlisting}[style=mgpy]
meta,scenario_1,scenario_1
hour,2026,2027
0,15.2,16.0
1,14.9,15.8
\end{lstlisting}

The loaded array has dimensions \code{(year, period, scenario)}.

\section{\path{resource_availability.csv}}

This file uses a three-row header:

\begin{itemize}[leftmargin=1.5em]
  \item row 1: scenario,
  \item row 2: year,
  \item row 3: resource.
\end{itemize}

The first column is \code{("meta","hour","")}. All scenario-year-resource combinations must exist. The file must contain exactly 8760 rows. Values are intended as capacity factors close to the interval $[0,1]$, although the current loader tolerates a broader numeric bound.

\begin{lstlisting}[style=mgpy]
meta,scenario_1,scenario_1
hour,2026,2026
,Solar,Wind
0,0.00,0.45
1,0.00,0.43
\end{lstlisting}

The loaded array has dimensions \code{(year, period, scenario, resource)}.

\section{\path{grid_import_price.csv} and \path{grid_export_price.csv}}

These files follow the same two-row scenario/year header used by \path{load\_demand.csv}. They are required only for on-grid projects, with \path{grid\_export\_price.csv} additionally conditional on export being enabled. Values are expressed in currency per kWh.

\section{\path{generator_efficiency_curve.csv}}

Required columns are:

\begin{itemize}[leftmargin=1.5em]
  \item \code{Relative Power Output [-]}
  \item \code{Efficiency [-]}
\end{itemize}

Current accepted behavior:

\begin{itemize}[leftmargin=1.5em]
  \item an explicit \code{0.0} row is optional,
  \item at least one positive-load point is required,
  \item the last relative-power point must be \code{1.0},
  \item positive-load points must be strictly increasing,
  \item efficiency values for positive-load points must be strictly positive.
\end{itemize}

The preferred interpretation is that \code{Efficiency [-]} is a normalized multiplier relative to \code{generator.technical.nominal\_efficiency\_full\_load}. Legacy absolute-efficiency curves are still accepted where supported by the parser.

\section{\path{battery_efficiency_curve.csv}}

Required columns are:

\begin{itemize}[leftmargin=1.5em]
  \item \code{relative\_power\_pu}
  \item \code{charge\_efficiency}
  \item \code{discharge\_efficiency}
\end{itemize}

The file must contain at least two rows, \code{relative\_power\_pu} must be strictly increasing in the interval $(0,1]$, and the last point must be \code{1.0}. The preferred interpretation is that the charge and discharge columns are normalized multipliers around the scalar efficiencies declared in \path{battery.yaml}.

\section{\path{battery_calendar_fade_curve.csv}}

Required columns are:

\begin{itemize}[leftmargin=1.5em]
  \item \code{soc\_pu}
  \item preferred: \code{calendar\_fade\_coefficient\_per\_year}
  \item legacy alias: \code{calendar\_fade\_coefficient\_per\_step}
\end{itemize}

The file must contain at least two rows. \code{soc\_pu} must lie in $[0,1]$. If the first point is above zero, the loader prepends a zero point automatically. The last \code{soc\_pu} point must be \code{1.0}. The coefficient is interpreted as a yearly calendar-ageing coefficient used together with \code{battery\_calendar\_time\_increment\_per\_year}.

\section{\path{grid_availability.csv}}

This file is derived rather than user-authored. It is regenerated from \path{grid.yaml} and the current Multi-Year sets. It uses the same two-row scenario/year header as the price files and represents binary hourly availability.

\chapter{Multi-Year Conditional Features}\label{sec:multi-year-conditional}

\section{On-grid mode}

Activated by \code{on\_grid = true}. Additional required inputs are:

\begin{itemize}[leftmargin=1.5em]
  \item \path{grid.yaml}
  \item \path{grid\_import\_price.csv}
  \item \path{grid\_export\_price.csv} if export is enabled
\end{itemize}

The derived file is \path{grid\_availability.csv}.

\section{Export mode}

Activated by \code{grid\_allow\_export = true}. This adds \path{grid\_export\_price.csv} and activates export-side variables in the optimization model.

\section{Capacity expansion}

Activated by \code{capacity\_expansion = true}. Renewable, battery, and generator templates then generate one \code{investment.by\_step} block per step, and the model tracks each commissioned cohort by \code{inv\_step}.

\section{Multi-scenario mode}

Activated by \code{multi\_scenario.enabled = true}. All hourly time-series CSVs must then contain every scenario-year combination. \code{fuel.cost.by\_scenario} and \code{grid.by\_scenario} must also cover every scenario label.

\section{Battery convex loss model}

Activated by \code{battery\_model.loss\_model = "convex\_loss\_epigraph"}. This requires \path{battery\_efficiency\_curve.csv} and switches the battery-loss representation from scalar efficiencies to a convex piecewise-linear surrogate based on DC-side relative power.

\section{Battery endogenous degradation}

Activated by convex battery loss together with either \code{cycle\_fade\_enabled} or \code{calendar\_fade\_enabled}. Additional battery inputs may include:

\begin{itemize}[leftmargin=1.5em]
  \item \code{initial\_soh},
  \item \code{end\_of\_life\_soh},
  \item \code{cycle\_lifetime\_to\_eol\_cycles},
  \item \path{battery\_calendar\_fade\_curve.csv},
  \item \code{calendar\_time\_increment\_per\_year}.
\end{itemize}

Cycle fade and calendar fade are tracked as yearly degradation state variables. Replacement timing still follows the declared battery calendar lifetime.

\section{Generator efficiency curve}

Activated by \code{generator\_model.efficiency\_model = "efficiency\_curve"}. This requires \path{generator\_efficiency\_curve.csv} and converts the user curve into an internal convex fuel-use surrogate.

\chapter{Multi-Year Consistency Checks}\label{sec:multi-year-checks}

Before solving, verify the following:

\begin{enumerate}[leftmargin=1.5em]
  \item \code{start\_year\_label} is an integer-like calendar year.
  \item If capacity expansion is enabled, \code{investment\_steps\_years} sums exactly to \code{time\_horizon\_years}.
  \item Scenario labels match in \path{formulation.json}, CSV headers, \code{fuel.cost.by\_scenario}, and \code{grid.by\_scenario}.
  \item Resource labels match between \code{renewables[].resource} and the third header row of \path{resource\_availability.csv}.
  \item \code{fuel.cost.by\_scenario.<scenario>.by\_year\_cost\_per\_unit\_fuel} contains one value for every modeled year.
  \item All hourly files contain exactly 8760 rows and the correct \code{meta/hour} sequence from 0 to 8759.
  \item If on-grid mode is active, \code{first\_year\_connection} uses an exact modeled year label whenever possible.
  \item If land constraints are active, every renewable entry provides a meaningful \code{specific\_area\_m2\_per\_kw}.
  \item If convex battery loss is active, \code{battery.technical.efficiency\_curve\_csv} points to an existing curve file.
  \item If battery calendar fade is enabled, \path{battery\_calendar\_fade\_curve.csv} is present and the calendar time increment uses the supported constant-per-year semantics.
  \item If generator curve mode is active, the generator curve ends at relative power \code{1.0}.
  \item Do not try to make technical parameters vary by investment step. The current Multi-Year backend assumes shared technical parameters across steps.
\end{enumerate}

\chapter{Multi-Year Minimal Examples}\label{sec:multi-year-examples}

\section{Example A: simplest off-grid Multi-Year case}

This is the smallest dynamic project that still exercises the Multi-Year machinery:

\begin{itemize}[leftmargin=1.5em]
  \item \code{core\_formulation = "dynamic"}
  \item \code{system\_type = "off\_grid"}
  \item \code{on\_grid = false}
  \item \code{capacity\_expansion = false}
  \item \code{start\_year\_label = 2026}
  \item \code{time\_horizon\_years = 10}
  \item \code{multi\_scenario.enabled = false}
  \item one renewable resource
  \item battery loss model = \code{constant\_efficiency}
  \item generator efficiency model = \code{constant\_efficiency}
\end{itemize}

Required files are \path{formulation.json}, \path{load\_demand.csv}, \path{resource\_availability.csv}, \path{renewables.yaml}, \path{battery.yaml}, and \path{generator.yaml}. In the single-step case, each component uses one \code{"1"} investment block.

\section{Example B: richer on-grid Multi-Year case}

This configuration exercises most dynamic features:

\begin{itemize}[leftmargin=1.5em]
  \item \code{core\_formulation = "dynamic"}
  \item \code{system\_type = "on\_grid"}
  \item \code{on\_grid = true}
  \item \code{grid\_allow\_export = true}
  \item \code{capacity\_expansion = true}
  \item \code{investment\_steps\_years = [5,5]}
  \item \code{start\_year\_label = 2026}
  \item \code{time\_horizon\_years = 10}
  \item optional multiple scenarios
  \item optional convex battery loss and endogenous degradation
  \item optional generator efficiency curve
\end{itemize}

Additional required files are \path{grid.yaml}, \path{grid\_import\_price.csv}, \path{grid\_export\_price.csv}, and any active advanced-curve files.

\part{Typical-Year Planning Mode}

\chapter{Typical-Year Overview}\label{sec:typical-overview}

Typical-Year is implemented as \code{core\_formulation = "steady\_state"}. It uses the same application pages and the same general project workflow defined in Part~I, but it collapses the dynamic structure into a representative-year planning problem.

Conceptually, Typical-Year is a reduced special case of the Multi-Year workflow described in Part~II:

\begin{itemize}[leftmargin=1.5em]
  \item the operating year is represented once and repeated conceptually,
  \item investment costs are annualized directly,
  \item the objective minimizes expected equivalent annual cost,
  \item scenarios may differ operationally while sharing the same design decisions.
\end{itemize}

\begin{guidenote}[Typical-Year simplification]
Typical-Year follows the same platform workflow, file layout, and component logic described for Multi-Year, but removes explicit year evolution, explicit investment-step dynamics, and battery-degradation state tracking across years.
\end{guidenote}

\chapter{Differences from Multi-Year}\label{sec:typical-differences}

This chapter lists only the differences from the reference Multi-Year workflow in Sections~\ref{sec:multi-year-init} to~\ref{sec:multi-year-examples}.

\section{Initialization differences}

When Typical-Year is selected, Project Setup writes:

\begin{itemize}[leftmargin=1.5em]
  \item \code{core\_formulation = "steady\_state"},
  \item \code{start\_year\_label = "typical\_year"}.
\end{itemize}

There is no dynamic horizon and no explicit Multi-Year cohort accounting in the optimization path. The representative year is indexed by the single label \code{typical\_year}.

\section{Investment-step simplification}

Typical-Year still uses the same YAML nesting concept, namely \code{investment.by\_step}, but only one step is relevant. Current Typical-Year templates keep the single-step key \code{base}. The Typical-Year parsers also accept a single non-\code{base} step block when only one block is present.

\section{CSV simplification}

Typical-Year uses the same CSV file families described in Chapter~\ref{sec:multi-year-ts}, with the following simplifications:

\begin{itemize}[leftmargin=1.5em]
  \item the year header is the single label \code{typical\_year},
  \item there is no explicit dynamic year expansion,
  \item fuel prices are scalar scenario values rather than yearly trajectories,
  \item grid connection timing by explicit year does not apply.
\end{itemize}

\section{YAML simplification}

Typical-Year reuses the same component files described in Chapter~\ref{sec:multi-year-inputs}, with the following simplifications:

\begin{itemize}[leftmargin=1.5em]
  \item investment data are single-step rather than step-dependent over a horizon,
  \item battery degradation inputs such as \code{initial\_soh}, calendar-fade curves, and endogenous year-to-year degradation states are not part of the current Typical-Year workflow,
  \item generator fuel data use scenario-based scalar fuel costs rather than yearly price lists,
  \item \code{first\_year\_connection} is not used because the representative-year formulation does not model explicit calendar years.
\end{itemize}

\section{Conditional-feature simplification}

Most conditional features remain conceptually the same as in Chapter~\ref{sec:multi-year-conditional}. The main formulation-specific simplifications are:

\begin{itemize}[leftmargin=1.5em]
  \item convex battery loss remains available,
  \item generator efficiency curves remain available,
  \item explicit battery endogenous degradation is not part of the current Typical-Year user workflow,
  \item capacity expansion and year-dependent cohort logic are not active in the optimization path.
\end{itemize}

\section{Consistency-check differences}

In addition to the shared file-consistency checks inherited from Chapter~\ref{sec:multi-year-checks}, the Typical-Year user should verify:

\begin{itemize}[leftmargin=1.5em]
  \item all time-series files use the single year header \code{typical\_year},
  \item only one investment block is provided in each component file,
  \item on-grid projects provide the required grid-price templates but do not attempt to declare dynamic first-connection years,
  \item dynamic battery-degradation inputs are not introduced into Typical-Year files.
\end{itemize}

\part{Appendices}

\appendix

\chapter{Glossary}\label{sec:appendix-glossary}

\begin{longtable}{p{0.22\textwidth}p{0.70\textwidth}}
\toprule
\textbf{Term} & \textbf{Meaning in current implementation} \\
\midrule
\endfirsthead
\toprule
\textbf{Term} & \textbf{Meaning in current implementation} \\
\midrule
\endhead
\code{year} & explicit modeled year label, usually a calendar year such as \code{2026}; in Typical-Year this collapses to \code{typical\_year} \\
\code{period} & hourly index within the representative operating year, \code{0,\dots,8759} \\
\code{scenario} & uncertainty branch with shared design decisions and different operating conditions \\
\code{inv\_step} & investment step and commissioning cohort index \\
\code{resource} & renewable-resource label used in sets and time-series headers \\
cohort / vintage & capacity commissioned in an investment step and then tracked through later years \\
shared-technology model & current Multi-Year interpretation in which investment terms may vary by step but technical parameters remain shared \\
capacity expansion & activation of multiple investment steps over a dynamic horizon \\
\bottomrule
\end{longtable}

\chapter{UI Mapping}\label{sec:appendix-ui}

\begin{longtable}{p{0.32\textwidth}p{0.60\textwidth}}
\toprule
\textbf{UI choice} & \textbf{Effect} \\
\midrule
\endfirsthead
\toprule
\textbf{UI choice} & \textbf{Effect} \\
\midrule
\endhead
Multi-Year formulation & writes \code{core\_formulation = "dynamic"} \\
Typical-Year formulation & writes \code{core\_formulation = "steady\_state"} \\
Start year and horizon & set the Multi-Year \code{year} coordinate and year headers \\
Capacity expansion enabled & generates multiple \code{investment.by\_step} blocks in Multi-Year \\
On-grid mode & generates \path{grid.yaml} and \path{grid\_import\_price.csv} \\
Export enabled & also generates \path{grid\_export\_price.csv} \\
Multi-scenario mode & replicates scenario labels across CSV headers and scenario-keyed YAML mappings \\
Battery loss model = \code{convex\_loss\_epigraph} & generates the battery efficiency-curve template and activates curve-based loader logic \\
Battery calendar fade enabled & generates the calendar-fade curve template and writes calendar-fade battery fields \\
Generator efficiency model = \code{efficiency\_curve} & generates the generator efficiency-curve template and writes its filename into \path{generator.yaml} \\
\bottomrule
\end{longtable}

\chapter{Source Basis}\label{sec:appendix-sources}

This unified guide was reconstructed from the current repository implementation, especially:

\begin{itemize}[leftmargin=1.5em]
  \item \path{pages/0\_Project\_Setup.py}
  \item \path{pages/1\_Data\_Audit\_and\_Visualization.py}
  \item \path{pages/3\_Optimization.py}
  \item \path{pages/4\_Results.py}
  \item \path{core/io/templates.py}
  \item \path{core/multi\_year\_model/sets.py}
  \item \path{core/multi\_year\_model/data.py}
  \item \path{core/data\_pipeline/battery\_loss\_model.py}
  \item \path{core/data\_pipeline/battery\_degradation\_model.py}
  \item \path{core/data\_pipeline/battery\_calendar\_fade\_model.py}
  \item \path{core/data\_pipeline/generator\_partial\_load\_model.py}
\end{itemize}

When code, template comments, and older prose diverged, this guide followed the current implementation and the current generated-template workflow.

\end{document}
